\section{Endpoint response of the repeated geodesics}
\label{sec:geodesic-jacobi}

Let \(\boldsymbol P=r^2\dd\boldsymbol{x}/\dd s\) be the conserved
transverse momentum per unit worldline mass. Allow it to vary near
zero while transporting the turning points and the contour. The
radial equation becomes
\begin{equation}
 Y^2=E^2+f(r)\left(1-\frac{\boldsymbol P^2}{r^2}\right).
 \label{eq:geodesic-transverse}
\end{equation}
Write \(\boldsymbol X\) for the spatial separation of the
endpoints. At \(\boldsymbol P=0\), the Jacobi map for each transverse
direction is given by a single radial integral,
\begin{equation}
 \left.\frac{\partial X_i}{\partial P_j}\right|_{E,\boldsymbol P=0}
   =\delta_{ij}J,\qquad
 J=\int_{\rm contour}\frac{\dd r}{r^2Y}.
 \label{eq:geodesic-jacobi-integral}
\end{equation}
Indeed \(\dd\boldsymbol{x}=\boldsymbol P\,\dd s/r^2\), and the
variation of \(r\) starts at order \(\boldsymbol P^2\). The
turning points are simple, so their contributions to the
linearization are regular in the square-root coordinate.
The time response is even in \(\boldsymbol P\), and at
\(\boldsymbol P=0\) the mixed time--space response vanishes.

Consider the contour of section~\ref{sec:optical-cycles} with \(n-1\)
optical cycles and insert \(k\geq0\) copies of the adjacent-root
cycle of section~\ref{sec:cycles}. The response is additive at
common energy,
\begin{equation}
 J_{n,k}=J_0+(n-1)(J_0-J_{\rm ext})+k\Pi_J,\qquad
 \Pi_J=\oint_{\mathcal C}\frac{\dd r}{r^2Y}.
 \label{eq:jacobi-additivity}
\end{equation}
Here \(J_0\) is for the reference path and \(J_{\rm ext}\) is twice
the integral from \(r_{\rm b}\) to infinity on the sheet used in
\eqref{eq:optical-outer-factor}. These integrals converge at
infinity. The integrand is regular at a horizon, so the horizon
detour does not affect \(J\).

\subsection{Transverse periods in the null limit}

For the reference path, let \(y=r/r_0\) and
\(\mu=(r_0/E)^2\). The exact integral is
\begin{equation}
 J_0=\frac{2}{Er_0}\int_1^\infty
 \frac{\dd y}{y^2\sqrt{1-y^{-m}+\mu(y^2-y^{-m})}}.
 \label{eq:jacobi-reference}
\end{equation}
At \(\mu=0\) the integral gives \(B(1/m,1/2)/m\).
The first correction comes from the exterior scaling
\(z=\sqrt{\mu}\,y\) in the difference between the integrals
at \(\mu\) and zero. Dividing the difference by \(\sqrt\mu\)
gives the limit
\begin{equation}
 \int_0^\infty\frac{(1+z^2)^{-1/2}-1}{z^2}\,\dd z=-1 .
\end{equation}
On \(1\leq y\leq2\) the difference contributes \(O(\mu)\).
For \(y\geq2\), the subtracted integrand in \(z\) is bounded
by a constant times \(\min(1,z^{-2})\).
Using \(r_0=E^{-2/m}[1+O(E^{-2q})]\) therefore gives
\begin{equation}
 J_0=c_d E^{q-2}-2E^{-2}+o(E^{-2}),\qquad
 c_d=\frac{2}{m}B(1/m,1/2).
 \label{eq:jacobi-reference-asymptotic}
\end{equation}
The exterior integral in \eqref{eq:jacobi-additivity} uses the same
leading integral as the exterior time in section~\ref{sec:optical-cycles},
with factor \(2/r_{\rm b}^{\,2}\). It equals
\(-2E^{-2}+O(E^{-d-2})\).

For \(d>4\), each cap of \(\Pi_J\) has the leading differential
\(-\ii r^{m/2-2}\dd r/\sqrt{1-E^2r^m}\). With
\(\alpha_J=1/2-1/m>0\), the transported phase on the next cap
is \(e^{2\pi\ii\alpha_J}\). Factoring the simple root as in
\eqref{eq:root-factor} gives the leading period
\begin{align}
 \Pi_J(E)&=
 -\frac{2\ii}{m}(e^{2\pi\ii\alpha_J}-1)B(\alpha_J,1/2)
       E^{q-2}+O(E^{-q-2})\notag\\
 &= c_d(1-e^{-2\pi\ii/m})E^{q-2}+O(E^{-q-2}).
 \label{eq:jacobi-local-period}
\end{align}
The second equality uses the gamma reflection identity. In AdS\(_5\),
the two cap integrals cancel by \(r\mapsto-r\), with the logarithm
transported through the upper half-plane. The arc at \(r=0\)
now has the leading differential \(-\ii\dd r/r\). Its contribution
on the half-turn is \(\pi\), doubled by the return on the other
sheet. Consequently \(\Pi_J=2\pi\) exactly. This is also the
\(m\to2\) limit of the coefficient in \eqref{eq:jacobi-local-period}.

Combining the optical and local periods yields
\begin{equation}
 J_{n,k}(E)
 =c_d\big[n+k(1-e^{-2\pi\ii/m})\big] E^{q-2}
       -2E^{-2}+o(E^{-2}).
 \label{eq:jacobi-combined}
\end{equation}
The \(E^{-2}\) term comes solely from the two endpoint legs in
this formula for fixed \(n,k\). In AdS\(_5\), using \(u=r^{-2}\)
reduces \(J_0\) to the integral of
\(\dd u/\sqrt{1+E^2u-u^2}\) from zero to its positive root.
The exterior integral has the same antiderivative with the
negative root as endpoint and the orientation from \(r_{\rm b}\)
toward infinity. These integrals give
\begin{equation}
 J_{n,k}^{\mathrm{AdS}_5}(E)
     =(n+2k)\pi-\arctan(2/E^2).
 \label{eq:jacobi-ads5}
\end{equation}

\begin{table}[t]
 \centering\small
 \begin{tabular}{ccc}
 \hline
 \(d\) & \(E\) & \(E^2\big\{J_{3,1}-c_d(4-e^{-2\pi\ii/m})E^{q-2}\big\}\) \\
 \hline
 4 & 8 & $-1.999349+0.000000i$ \\
4 & 128 & $-2.000000+0.000000i$ \\
5 & 8 & $-1.898401-0.019149i$ \\
5 & 128 & $-1.999005-0.000191i$ \\
6 & 8 & $-1.840734-0.039607i$ \\
6 & 128 & $-1.997518-0.000621i$ \\
8 & 8 & $-1.782179-0.053899i$ \\
8 & 128 & $-1.994597-0.001337i$ \\
\hline

 \end{tabular}
 \caption{Transverse response for three optical visits and one
 additional adjacent-root cycle. The entries approach \(-2\).}
 \label{tab:geodesic-jacobi}
\end{table}

For table~\ref{tab:geodesic-jacobi} the optical part is integrated
along the caps, horizon detour and arc used in section~\ref{sec:optical-cycles}.
The direct optical integral and the reference-minus-exterior result
agree to absolute error below \(10^{-12}\). The AdS\(_5\) values
also reproduce the exact reference integral and both cycle periods.

\subsection{Endpoint Hessian and Gaussian factor}

The endpoint variation of the length is
\(\dd L=-E\dd T+\boldsymbol P\cdot\dd\boldsymbol X\).
At \(\boldsymbol X=0\) this yields
\begin{equation}
 \partial_T^2 L=-(T'_{n,k})^{-1},\qquad
 \partial_T\partial_{X_i}L=0,\qquad
 \partial_{X_i}\partial_{X_j}L=\delta_{ij}J_{n,k}^{-1},
 \label{eq:jacobi-hessian}
\end{equation}
where \(T_{n,k}\) includes the two types of cycle.
The leading transverse coefficient has positive real part
for \(n\geq1\) and \(k\geq0\), so the endpoint map is invertible
at sufficiently large finite \(E\).
In particular, the first spatial correction at fixed
\(T\) is \(\boldsymbol X^2/(2J_{n,k})\). In AdS\(_5\) the
adjacent-root cycle has zero radial length and time periods
at every \(E\), while \eqref{eq:jacobi-ads5} gives distinct
quadratic responses for different \(k\).

The reduced action \(W=L+ET-\boldsymbol P\cdot\boldsymbol X\)
has time block \(T'_{n,k}\) and spatial block
\(-J_{n,k}\mathbf 1_{d-1}\) in its Hessian.
The endpoint Gaussian factor on the associated saddle is
\(\mathcal D_{n,k}=(T'_{n,k}J_{n,k}^{d-1})^{-1/2}\), up to
the phase of the Gaussian contour. For fixed \(n,k\),
the endpoint time tails give \(T'_{n,k}\sim 2/E^2\).
At a common small deficit \(\delta\) from the respective optical
times, each energy satisfies \(E\sim2/\delta\).
The ratio of endpoint Gaussian magnitudes is therefore
\begin{equation}
 \left|\frac{\mathcal D_{n,k}}{\mathcal D_{1,0}}\right|
 \longrightarrow
 \big|n+k(1-e^{-2\pi\ii/m})\big|^{-(d-1)/2}.
 \label{eq:jacobi-ratio}
\end{equation}
For the pure optical family the factors are \(n^{-(d-1)/2}\).
The Gaussian approximation in transverse momentum requires a
width \((M|J_{n,k}|)^{-1/2}\) small compared with \(r_0\).
For fixed \(n,k\) this means \(ME^{-q}\gg1\). The endpoint factor
describes that semiclassical regime. The contribution to a
correlator additionally requires the weight of each complex
saddle on its integration contour.

\subsection{Finite transverse momentum for the Lovelock lapse}
\label{sec:finite-transverse-geodesic}

The lapse \(f=r^2-r^{-2}\), which also occurs in the Lovelock
family of section~\ref{sec:lovelock}, permits a resummation of the
transverse response. Write \(P^2=\boldsymbol P^2\) and use
\(u=E^2r^2\), \(\eta=E^2P^2\). The exact equation is
\begin{equation}
 \frac{Y^2}{E^2}=1-\frac1u+\frac{\eta}{u^2}
       +E^{-4}(u-\eta).
\end{equation}
At fixed \(0\leq\eta<1/4\), the inner turn of the optical cycle
tends to \(u=a_+(\eta)\), with
\(a_\pm(\eta)=(1\pm\sqrt{1-4\eta})/2\).
The integral \(\int \dd r/(r^2Y)\) on the doubled positive ray
converges to the elliptic period
\begin{align}
 \mathcal J(\eta)
 &=\int_{a_+}^{\infty}
     \frac{\dd u}{\sqrt{u(u-a_+)(u-a_-)}} \notag\\
 &=\frac{2}{\sqrt{a_+}}\,\mathbb K(a_-/a_+),\qquad
 \mathbb K(m)=\int_0^{\pi/2}
      \frac{\dd\theta}{\sqrt{1-m\sin^2\theta}} .
 \label{eq:jacobi-elliptic}
\end{align}
The substitution \(u=a_+/\cos^2\theta\) removes the endpoint
square root and gives the second line. It also gives convergence
uniformly on compact \(\eta\) intervals before the two inner
roots coalesce. The doubled tail from the imaginary outer root
is \(O(E^{-2})\), since \(r_{\rm b}\sim\ii E\). Thus both the
reference path and the optical cycle have limiting transverse
period \(\mathcal J(\eta)\), and the sequence with \(n\) inner
turns has displacement \(n\boldsymbol P\,\mathcal J(\eta)+o(P)\).
At \(\eta=0\), \(\mathcal J(0)=\pi\) reproduces the linear
response. As the inner roots coalesce at \(\eta=1/4\), the
period diverges logarithmically \cite{NIST:DLMF}.

Let \(W_{\rm opt}=\oint Y\,\dd r/f\) on one optical cycle.
Differentiation at fixed \(E\) gives
\(\partial_{P^2}W_{\rm opt}=-\tfrac12\oint\dd r/(r^2Y)\).
Simple turning points give zero endpoint terms. Integrating
from zero transverse momentum therefore yields
\begin{align}
 W_{\rm opt}(E,P^2)-W_{\rm opt}(E,0)
   &=-\frac{\mathcal A(\eta)}{2E^2}+o(E^{-2}),\qquad
    \mathcal A(\eta)=\int_0^\eta\mathcal J(s)\,\dd s,\notag\\
 \mathcal A(\eta)&=\pi\big(\eta+\tfrac38\eta^2+O(\eta^3)\big).
 \label{eq:jacobi-elliptic-action}
\end{align}
Each additional optical cycle adds this same action shift.
The reference path contributes one copy to leading order, so
the shift for \(n\) inner visits has coefficient \(n\).

\begin{table}[t]
 \centering\small
 \begin{tabular}{ccccc}
 \hline
 \(\eta\) & \(\mathcal J(\eta)\) & \(E=8\) & \(E=32\) & \(E=128\) \\
 \hline
 0.0625 & 3.3132763 & 3.04e-07 & 1.19e-09 & 4.63e-12 \\
0.16 & 3.7694524 & 5.98e-06 & 2.34e-08 & 9.13e-11 \\
0.24 & 5.2387499 & 1.59e-04 & 6.22e-07 & 2.43e-09 \\
\hline

 \end{tabular}
 \caption{Relative errors of the finite-energy optical period
 compared with the elliptic limit in \eqref{eq:jacobi-elliptic}.}
 \label{tab:geodesic-momentum}
\end{table}

For table~\ref{tab:geodesic-momentum}, we compute the two
finite-energy turning points on the optical cycle and subtract
the imaginary-ray tail from the positive-ray integral.
Section~\ref{sec:lovelock-crossover} uses the analytic continuation
from \(\eta=0\) to compare the nonlinear action with the scalar
connection coefficient.
